\documentclass[11pt,a4paper]{article}
\usepackage[utf8]{inputenc}
\usepackage[T1]{fontenc}
\usepackage{xeCJK} 
\usepackage{amsmath}
\usepackage{graphicx}
\usepackage{geometry}
\usepackage{booktabs}
\usepackage{listings}
\usepackage{xcolor}
\usepackage{hyperref}

\geometry{left=2.5cm,right=2.5cm,top=2.5cm,bottom=2.5cm}

\title{CSE5023：深度学习前沿进展作业 2\\ Transformer 架构及其应用实验报告}
\author{段伟强 (Weiqiang Duan) \\ 学号：12533582}
\date{2025年4月26日}

\begin{document}

\maketitle

\section{第一部分：机器翻译任务 (50分)}

\subsection{对 Transformer 架构的思考与理解}

针对作业 1.1 节提出的核心问题，本人的理解如下：

\paragraph{1. 编码器（Encoder）与解码器（Decoder）在结构设计和功能定位上的差异：}
编码器（Encoder）采用双向注意力机制，其功能定位是“特征提取器”，旨在将源语言序列映射为高维语义表示。解码器（Decoder）则引入了交叉注意力层，用于对齐源语言特征；其功能定位是“生成器”，通过自回归方式产生目标语言。

\paragraph{2. 注意力掩码（Attention Mask）的实现与核心动机：}
在本实验中，编码器使用 Padding Mask 屏蔽填充位。解码器则额外使用了由下三角矩阵生成的自回归掩码（Causal Mask）。其核心动机是维护因果性，防止模型在预测当前词时“窥探”未来的真实标签，确保训练与推理的一致性。

\paragraph{3. 解码器在训练阶段和推理阶段的工作流程差异：}
训练阶段采用 Teacher Forcing，解码器可以并行接收完整的目标序列并计算损失；推理阶段则是逐词生成的自回归过程，上一步的输出作为下一步的输入，属于串行计算。

\subsection{模型训练与评估结果}

\paragraph{训练过程：} 本模型在全量数据集上训练了 50 轮。平均损失值（Loss）从初始的 5.4 稳步下降至 \textbf{0.3079}。训练曲线显示模型收敛良好，未出现明显的梯度消失或爆炸。

\paragraph{BLEU 定量指标：}
在 1817 句独立测试集上的最终评估结果如下：
\begin{table}[h]
\centering
\begin{tabular}{ll}
\toprule
指标项目 & 测量数值 \\
\midrule
\textbf{最终测试集 BLEU 分数} & \textbf{0.2121} (21.21\%) \\
1-gram Precision & 0.5699 \\
4-gram Precision & 0.0958 \\
Brevity Penalty (BP) & 0.9889 \\
Length Ratio & 0.9889 \\
\bottomrule
\end{tabular}
\end{table}

\subsection{翻译输出示例与性能分析}

以下是模型在测试集上的具体表现分析：

\begin{enumerate}
    \item \textbf{语义对齐示例}：
    \textit{Ref: 她在浴室。$\rightarrow$ Pred: 她在洗澡。} \\
    分析：虽然词汇不完全一致，但模型准确捕捉到了场景的语义关联，展示了较强的泛化能力。
    
    \item \textbf{复杂逻辑理解}：
    \textit{Ref: 他聪明到不会娶她。$\rightarrow$ Pred: 他不会傻到去娶她。} \\
    分析：这是一个非常成功的案例。模型不仅理解了否定逻辑，还使用了“傻”这个词与原文的“聪明”形成逻辑对冲，说明模型学习到了深层意图。
    
    \item \textbf{典型失效案例（极性反转）}：
    \textit{Ref: 这是我最差劲的电影。$\rightarrow$ Pred: 这是我见过最好的电影响。} \\
    分析：模型将“最差”识别成了“最好”。这揭示了 NMT 模型在处理情感极性相反但上下文语境相似的形容词时，仍存在识别难度。
    
    \item \textbf{幻觉现象分析}：
    \textit{Ref: 人是唯一会使用火的动物。$\rightarrow$ Pred: 人类是唯一能够彼此交动物的动物。} \\
    分析：当面对复杂的陈述句时，模型出现了语序混乱和词汇堆砌，这表明在有限的训练数据下，模型对复杂常识性句子的拟合仍有待加强。
\end{enumerate}

\subsection{预热策略与学习率调优}

\paragraph{实验动机与策略原理：}
为了进一步提升翻译质量，本实验引入了带预热的余弦退火策略（Cosine Annealing with Warm-up），替换基线实验中的固定学习率。该策略的核心机制为：在预热阶段，学习率从零线性上升至峰值；预热结束后，按半余弦曲线从峰值平滑下降至零。公式如下：
\[
\mathrm{lr}(\text{step}) =
\begin{cases}
\mathrm{peak\_lr} \cdot \dfrac{\text{step}}{\text{warmup\_steps}}, & \text{step} < \text{warmup\_steps} \\[6pt]
\mathrm{peak\_lr} \cdot 0.5 \left(1 + \cos\left(\pi \cdot \dfrac{\text{step} - \text{warmup\_steps}}{\text{total\_steps} - \text{warmup\_steps}}\right)\right), & \text{otherwise}
\end{cases}
\]

\paragraph{实验设计：}
设计了 3 组余弦退火实验与基线进行对比，每组均训练 50 个 epoch，通过验证集 BLEU 分数快速筛选最优参数，再用测试集进行最终评估：

\begin{table}[h]
\centering
\begin{tabular}{lccc}
\toprule
实验名称 & 学习率策略 & 峰值学习率 & 预热轮次 \\
\midrule
基线（固定学习率） & 固定 & $1\times 10^{-4}$ & 无 \\
实验A（余弦退火） & Cosine Warm-up & $1\times 10^{-4}$ & 5 \\
实验C（余弦退火） & Cosine Warm-up & $2\times 10^{-4}$ & 3 \\
\textbf{实验B} \textbf{（余弦退火）} & \textbf{Cosine Warm-up} & $\mathbf{3\times 10^{-4}}$ & \textbf{5} \\
\bottomrule
\end{tabular}
\end{table}

\paragraph{验证集 BLEU 动态对比：}
在训练过程中，每 10 个 epoch 在验证集上计算一次 BLEU 分数，结果如下：

\begin{table}[h]
\centering
\begin{tabular}{lcccc}
\toprule
训练轮次 & 实验A BLEU & 实验C BLEU & 实验B BLEU & 优胜方 \\
\midrule
Epoch 10 & -- & 0.1521 & 0.1736 & 实验B \\
Epoch 20 & -- & 0.2063 & \textbf{0.2268} & 实验B \\
Epoch 30 & -- & 0.2354 & \textbf{0.2460} & 实验B \\
Epoch 40 & 0.2001 & 0.2427 & \textbf{0.2556} & \textbf{实验B} $\uparrow$ \\
\bottomrule
\end{tabular}
\end{table}

实验A（peak\_lr=$1\times 10^{-4}$, warmup=5）的训练损失在后期陷入平台期（最终 Loss=0.914），验证 BLEU 仅达 0.2001，甚至低于基线水平。分析其原因为：峰值学习率与基线相同，但预热期过长（5 个 epoch），导致模型刚达到峰值学习率便开始衰减，有效学习时间严重不足，后期学习率趋近于零，模型基本停止收敛。

实验C（peak\_lr=$2\times 10^{-4}$, warmup=3）表现显著提升：训练损失持续下降至 \textbf{0.227}，验证 BLEU 在 epoch 40 达到 0.2427，相比实验A提升 21.3\%。

实验B（peak\_lr=$3\times 10^{-4}$, warmup=5）在所有阶段均领先：最终训练 Loss 低至 \textbf{0.084}（实验C的 1/3），验证 BLEU 在 epoch 40 达到 \textbf{0.2556}，相比实验C再提升 5.3\%。更大的峰值学习率配合充足的预热轮次，使得模型在整个训练周期内保持了更高的有效学习率，加速了收敛。

\paragraph{测试集最终评估：}
选取最优参数（实验B）在独立测试集上进行最终评估，与基线和实验C对比如下：

\begin{table}[h]
\centering
\begin{tabular}{llll}
\toprule
评估指标 & 基线（固定学习率） & 实验C（余弦退火） & \textbf{实验B} \textbf{（余弦退火）} \\
\midrule
\textbf{测试集 BLEU} & 0.2121 & 0.2268 & \textbf{0.2323} \\
最终训练 Loss & 0.308 & 0.227 & \textbf{0.084} \\
1-gram Precision & 0.5699 & 0.5840 & \textbf{0.5962} \\
4-gram Precision & 0.0958 & 0.1064 & \textbf{0.1111} \\
Brevity Penalty & 0.9889 & 0.9835 & 0.9672 \\
\bottomrule
\end{tabular}
\end{table}

实验B在测试集上取得了 \textbf{0.2323} 的 BLEU 分数，相比基线的 0.2121 提升了 \textbf{2.02 个百分点}（相对提升 9.5\%），相比实验C的 0.2268 也提升了 0.55 个百分点。各项 n-gram 精确率均为三组最高，其中 4-gram 从基线的 0.0958 提升至 0.1111（+16.0\%），说明更大峰值学习率的余弦退火策略显著改善了长短语的生成质量。

翻译质量直观对比（以实验B的最优结果与基线对照）：

\begin{itemize}
    \item \textit{Source: He hoped he could succeed, but he didn't.}
    \begin{itemize}
        \item 基线：他希望他们依然就没成功。（语义混乱）
        \item 实验B：他希望所有的成功，但他没有办法。（语义完整）
    \end{itemize}
    \item \textit{Source: Keep trying.}
    \begin{itemize}
        \item 基线：继续喝。（语义错误）
        \item 实验B：试试试。（接近原意）
    \end{itemize}
    \item \textit{Source: Take it.}
    \begin{itemize}
        \item 基线：它开它。（语法混乱）
        \item 实验B：把它拿去。（语法正确）
    \end{itemize}
\end{itemize}

\paragraph{实验结论：}
（1）余弦退火策略要发挥效果，峰值学习率必须适当高于固定学习率基线，否则预热期的开销将导致有效学习时间不足；（2）warmup 轮次不宜过短，当峰值学习率提高到 $3\times 10^{-4}$ 时，5 个 epoch 的预热能使学习率平稳过渡；（3）实验B（peak\_lr=$3\times 10^{-4}$, warmup=5）取得了最优结果，测试集 BLEU 达到 \textbf{0.2323}，相较于基线提升 \textbf{9.5\%}，充分验证了 Cosine Annealing with Warm-up 在机器翻译任务中的有效性；（4）三组实验呈现明显的单调递增趋势（1e-4 $\rightarrow$ 2e-4 $\rightarrow$ 3e-4），说明在 50 epoch 的余弦周期下，模型仍可承受更高的峰值学习率，更大峰值学习率的探索空间值得进一步研究。

\subsection{超参数消融实验}

\paragraph{实验设计：}
为了探究模型架构超参数对翻译性能的影响，选取实验B的最优配置作为基线（n\_layer=6, h\_num=8, d\_model=256, d\_ff=1024），保持训练策略（peak\_lr=$3\times 10^{-4}$, warmup=5, 50 epoch）不变，采用控制变量法逐一消融以下三个架构参数：

\begin{table}[h]
\centering
\begin{tabular}{lcccc}
\toprule
实验名 & n\_layer & h\_num & d\_model & d\_ff \\
\midrule
基线（实验B） & 6 & 8 & 256 & 1024 \\
消融1：减少层数 & 3 & 8 & 256 & 1024 \\
消融2：增加注意力头 & 6 & 16 & 256 & 1024 \\
消融3：降低嵌入维度 & 6 & 8 & 128 & 512 \\
\bottomrule
\end{tabular}
\end{table}

\paragraph{验证集 BLEU 动态对比：}

\begin{table}[h]
\centering
\begin{tabular}{lcccc}
\toprule
训练轮次 & 基线（n\_layer=6, h\_num=8, d\_model=256） & 消融1（n\_layer=3） & 消融2（h\_num=16） & 消融3（d\_model=128） \\
\midrule
Epoch 10 & 0.1736 & 0.1874 & 0.1539 & 0.0754 \\
Epoch 20 & 0.2268 & \textbf{0.2346} & 0.2211 & 0.1704 \\
Epoch 30 & 0.2460 & \textbf{0.2478} & 0.2297 & 0.1894 \\
Epoch 40 & 0.2556 & \textbf{0.2620} & 0.2297 & 0.2096 \\
最终 Loss & 0.084 & 0.121 & 0.078 & 0.729 \\
\bottomrule
\end{tabular}
\end{table}

\paragraph{结果分析：}

\textbf{消融1（n\_layer=3）：}令人意外的是，将 Transformer 层数从 6 层减半为 3 层后，验证集 BLEU 在各个阶段均略优于基线，epoch 40 达到 0.2620（基线 0.2556）。这表明在当前数据规模和任务复杂度下，3 层编码器-解码器已提供足够的表示容量，更多层数反而可能导致过拟合或优化困难。同时，3 层模型的训练速度（约 37 it/s）相比基线（约 20 it/s）提升了近 85\%，计算效率显著提高。

\textbf{消融2（h\_num=16）：}注意力头数从 8 增加到 16 后，验证 BLEU 从基线的 0.2556 下降至 0.2297（-10.1\%）。虽然最终 Loss（0.078）略低于基线（0.084），但 BLEU 表现不佳，说明更多的注意力头数并未带来更好的泛化能力。可能原因包括：在 d\_model=256 固定的前提下，头数翻倍意味着每个头的维度减半（256/16=16），过小的子空间限制了注意力的表达能力。

\textbf{消融3（d\_model=128）：}嵌入维度从 256 减半至 128 后，模型性能全面下降：epoch 40 BLEU 仅为 0.2096（-18.0\%），且最终 Loss（0.729）显著高于基线的 0.084。d\_model 决定了整个模型的表示容量，降低嵌入维度后模型无法充分学习源语言到目标语言的复杂映射关系，损失曲线在后期陷入平台期。此外，d\_ff=512 限制了前馈网络的非线性变换能力，进一步加剧了性能损失。

\paragraph{实验结论：}
（1）Transformer 层数对性能的影响呈现 U 型曲线：3 层已能取得与 6 层相当甚至略优的翻译质量，在资源受限场景下优先选择减少层数而非其他参数；（2）增加注意力头数需同步增加嵌入维度以保证每个头的子空间大小，否则将损害性能；（3）嵌入维度（d\_model）是影响模型容量最关键的超参数，盲目减半会导致严重的性能瓶颈；（4）计算效率方面，减少层数带来最大的加速收益（+85\%），而调整头数和嵌入维度的影响相对较小。

\section{第二部分：情感分析任务 (30分)}
(该部分实验正在进行中。初步计划加载翻译任务的编码器参数，通过冻结部分卷积层实现微调，预期能利用翻译任务学到的语义特征加速分类任务的收敛。)

\section{第三部分：视觉 Transformer (20分)}
(该部分计划对比 ViT 与 ResNet50 在图像识别上的表现。Transformer 架构的全局注意力机制在处理图像分块的关联性上具有独特优势。)

\end{document}
